\section{下游应用不是单一 AI 数据中心}
GPU 数量增加并不会线性等同于光通信全链条需求，真正的变量是集群网络拓扑、东西向流量、交换层级、冗余配置和光电转换边界。NVIDIA 与 Coherent/Lumentum 的合作说明系统厂已经把光互连视为集群扩展的关键瓶颈，而 Broadcom 102.4Tbps 交换芯片路线说明 1.6T 端口将成为下一代交换网络的自然配套。但完整光通信还包括 DCI/相干传输、运营商骨干、FTTH、企业园区、工业和汽车光互连，不同应用的订单节奏、利润弹性和估值倍数差异很大。

\begin{exhibitbox}[表：需求传导路径]
\centering
\small
\begin{tabularx}{\exhibitboxwidth}{>{\bfseries\raggedright\arraybackslash}p{2.0cm} >{\raggedright\arraybackslash\sloppy\hspace{0pt}}X >{\raggedright\arraybackslash}p{3.0cm}}
\toprule
\textbf{环节} & \textbf{传导逻辑} & \textbf{监测指标} \\
\midrule
GPU/ASIC 集群 & 参数同步和推理流量提升，交换层级增加 & NVIDIA/AMD/云厂商 capex 与集群出货节奏 \\
交换 ASIC & 51.2T 到 102.4T 切换推动端口速率升级 & Broadcom/Marvell 交换芯片发布时间表 \\
光模块 & 800G 继续放量，1.6T 开始资格认证 & 模块厂订单、ASP、毛利率、存货和现金流 \\
上游器件 & 激光器、光引擎、精密耦合和测试价值量提升 & 源杰/天孚/太辰光 Q2-Q3 收入和毛利率 \\
\bottomrule
\end{tabularx}
\end{exhibitbox}
\begin{exhibitbox}[表：下游应用与订单传导]
\centering\scriptsize
\renewcommand{\arraystretch}{1.18}
\begin{tabularx}{\exhibitboxwidth}{>{\bfseries\raggedright\arraybackslash}p{1.8cm} >{\raggedright\arraybackslash\sloppy\hspace{0pt}}p{3.0cm} >{\raggedright\arraybackslash\sloppy\hspace{0pt}}p{3.0cm} >{\raggedright\arraybackslash\sloppy\hspace{0pt}}X}
\toprule
\textbf{应用} & \textbf{拉动环节} & \textbf{A 股映射} & \textbf{投资节奏} \\
\midrule
AI 数据中心 & 800G/1.6T 光模块、交换 ASIC、光引擎、DSP、EML/CW 光源 & 中际旭创、新易盛、天孚通信、源杰科技、长光华芯 & 短周期最强，订单和毛利率必须季季验证 \\
DCI/相干传输 & 相干模块、ROADM/OTN、WDM、长距光纤 & 光迅科技、烽火通信、中兴通讯、长飞光纤 & 从云数据中心外溢到城域/长距互连，节奏慢于 AI 集群 \\
运营商骨干/5G/6G 承载 & OTN、PON、路由交换、光纤光缆、接入设备 & 中兴通讯、烽火通信、亨通光电、中天科技、长飞光纤 & 由运营商 capex 驱动，盈利弹性低但底盘稳定 \\
FTTH/企业园区 & ODN、PON、连接器、配线、室内外光缆 & 长飞光纤、亨通光电、中天科技、通鼎互联、永鼎股份、太辰光 & 价格竞争更强，适合作为现金流和周期底仓验证 \\
工业/汽车/边缘互连 & 短距光模块、连接器、传感与工业光器件 & 光库科技、太辰光、华工科技 & 商业化节奏分散，暂不作为主估值锚 \\
\bottomrule
\end{tabularx}
\end{exhibitbox}

\section{下游应用深挖：不同场景对应不同估值节奏}
\subsection{AI 数据中心}
AI 数据中心是本轮光通信最强的增量来源，但它拉动的是一组具体环节：交换 ASIC 端口升级、NIC/网卡带宽、800G/1.6T 光模块、DSP/TIA/driver、EML/CW 激光器、精密耦合和高速测试。投资上最直接的映射是中际旭创、新易盛，其次是天孚通信、源杰科技、长光华芯、仕佳光子、光库科技等上游芯片/器件。需要跟踪的不是“AI capex 是否高”，而是这些 capex 是否转化为可见订单、稳定 ASP 和毛利率。
\subsection{DCI 与相干传输}
DCI 和相干传输是 AI 数据中心外溢后的第二层需求。大型数据中心之间需要更高带宽、低功耗和长距离互连，相干模块、ROADM/OTN、WDM、光纤和网络设备都会受益。它的节奏通常慢于 AI 集群内部互连，但订单更偏系统项目，适合关注光迅科技、烽火通信、中兴通讯以及光纤光缆链条的项目交付。
\subsection{运营商骨干与 5G/6G 承载}
运营商骨干和承载网是光通信传统大盘。它不会像 AI 模块一样产生爆发式弹性，但决定行业底盘。中兴通讯、烽火通信、长飞光纤、亨通光电、中天科技、通鼎互联、永鼎股份受运营商 capex、招标份额、设备利润率、光缆价格和通信线缆项目影响更大。估值应看订单周期和现金流，不应简单跟随 AI 模块情绪。
\subsection{FTTH、企业园区和边缘节点}
FTTH、企业园区和边缘节点需求更分散，对 ODN、PON、连接器、室内外光缆、通信线缆和接入设备形成稳定拉动。该场景的竞争更充分、价格压力更明显，但可提供长周期收入底盘。太辰光、长飞光纤、亨通光电、中天科技、通鼎互联、永鼎股份和部分网络设备公司需要用现金流和项目利润率验证。
\subsection{工业、汽车和新型光互连}
工业和汽车光互连目前仍处于分散商业化阶段，可能在传感、车载高速互连、工业控制和边缘计算中形成长期需求，但 2026 年不宜作为主要估值分母。光库科技、华工科技和部分精密器件公司具备相关选项，本报告只把它作为长期上修因子，等待收入占比和客户结构更清晰后再提升权重。

\section{AI 快周期与运营商慢周期的估值差异}
AI 数据中心是 2025-2026 年最强的价格驱动因素，因为海外云厂商 capex、800G/1.6T 模块订单和毛利率能在几个季度内反映到利润表。运营商骨干、FTTH、DCI 和企业园区则更偏项目制和资本开支周期，需求没有消失，但估值应更多锚定现金流、订单能见度和资产周转，而不是套用 AI 模块的高增长倍数。材料和设备环节进一步滞后，它们可能先体现在产能和良率瓶颈，再进入收入确认。

\section{为什么不能把行业高景气直接等同于个股买入}
光通信是高景气行业，但股价买的是“高景气减去已经支付的估值”。2026-06-26 当天多个光通信标的下跌 5--10\%，但业务模型匹配后的基准估值仍显示不少公司当前价已经提前支付了较多成长和主题期权。只有当下一季度利润、收入质量、现金流和订单能见度继续超过估值隐含要求，主题才能继续转化为超额收益；否则股价会先通过横盘或下跌消化叙事。对全产业链而言，最重要的是把“短期 AI 模块兑现”“上游芯片/设备期权”“光纤光缆慢周期底盘”“网络设备应用层”分开看。
